\section{Application Developer Guide}
\label{sec:ApplicationDeveloper}

This section is for developers who want to use the library to give an
application Raft cluster behavior.

\subsection{What an application must provide}

An application supplies four things:

\begin{enumerate}
\item A command encoding for writes.
\item A deterministic state machine that applies committed commands.
\item A snapshot encoding for the application state.
\item Optional query handling for local reads after Raft has established read safety.
\end{enumerate}

The application does not implement leader election, quorum calculation, log
replication or membership transitions.

\subsection{Command design}

Commands are opaque byte payloads to Raft. They should be designed as a stable
application protocol:

\begin{itemize}
\item Use explicit command types rather than positional assumptions.
\item Include stable command identifiers when client retries must be idempotent.
\item Keep commands deterministic: all replicas must produce the same state from
the same ordered log.
\item Avoid embedding local time, random values or local host state unless those
values are part of the command chosen before replication.
\item Treat command formats as versioned public contracts once multiple deployed
nodes may run different software versions.
\end{itemize}

For reference-data style applications, typical commands are
\texttt{UPSERT\_PRODUCT}, \texttt{REMOVE\_PRODUCT},
\texttt{UPSERT\_VARIANT} and \texttt{REMOVE\_VARIANT}. For a key-value
application, commands may be \texttt{PUT}, \texttt{CAS}, \texttt{DELETE} and
\texttt{CLEAR}.

\subsection{State-machine contract}

The state machine is called only for committed entries. It should therefore be
written as pure replicated state mutation, not as a client-facing validation
service. Validation that decides whether a write is accepted before entering
Raft belongs in admission, authentication or authorization policy. Validation
that is part of deterministic domain semantics may happen inside the state
machine.

\begin{lstlisting}[language=Java,caption={Java state-machine shape}]
public interface SnapshotStateMachine {
    void apply(long term, byte[] command);
    byte[] snapshot();
    void restore(byte[] snapshotData);
}

public interface QueryableStateMachine extends SnapshotStateMachine {
    byte[] query(byte[] request);
}
\end{lstlisting}

The C++ interface carries the same idea using C++ types:

\begin{lstlisting}[language=C++,caption={C++ state-machine shape}]
class ApplicationStateMachine {
public:
    virtual std::string apply(std::int64_t index,
                              std::int64_t term,
                              std::string_view command) = 0;
    virtual std::string snapshot() const = 0;
    virtual std::string query(std::string_view request) const = 0;
    virtual void restore(std::string_view snapshot) = 0;
};
\end{lstlisting}

\subsection{Snapshots}

Application snapshots contain application state only. Raft wraps that payload
with Raft metadata, including the cluster configuration at the snapshot boundary.
This distinction matters because a snapshot restore must restore both:

\begin{itemize}
\item the application's lookup state,
\item the Raft membership and log boundary that make the snapshot safe.
\end{itemize}

Application developers should not store cluster configuration inside their
domain snapshot unless the domain itself has a reason to display or audit it.
Cluster configuration is Raft machinery state.

\subsection{Reads and read leases}

Applications may expose read queries through a queryable state machine. The
application should treat such reads as local state inspection. It should not try
to contact peers or decide whether the local node is a safe reader.

The runtime decides whether a read may be served:

\begin{itemize}
\item Followers redirect ordinary linearizable reads to the known leader.
\item A leader may serve a read if its quorum-backed read lease is fresh.
\item If the lease is stale, the leader must refresh a quorum barrier before serving.
\item If the barrier cannot be refreshed, the response should ask the client to retry.
\end{itemize}

\subsection{Client-facing policies}

The runtime has policy hooks around client commands:

\begin{description}
\item[Admission] Should this local node admit the write at all?
\item[Authentication] Who is making the request?
\item[Authorization] Is the authenticated principal allowed to perform this write?
\end{description}

These hooks are application-facing. They are not Raft safety hooks. A policy may
reject a write for domain or security reasons, but it must not claim a command is
committed. Only Raft commit and state-machine application can do that.

\subsection{Cluster configuration storage}

Applications do not need a separate domain table or object for cluster
configuration. Configuration is replicated through Raft internal log entries and
is captured in Raft snapshot metadata. The library persists it through the same
Raft persistence path as terms, votes and log state.

The important application responsibility is operational: the deployment must
provide durable storage for the Raft node. If the node loses its Raft storage it
may lose its term, vote, log and membership history. That is a cluster
operational failure, not a domain state-machine feature.
